% !TEX root = main.tex
\section*{Current Progress}
The implemented pipeline covers the full end-to-end path from a Chiron source program to a property verdict, encompassing state extraction, invariant generation, rule encoding, and property checking.

\subsection*{Literature Review}

\begin{itemize}[leftmargin=1.5em]
    \item \textbf{CHCs as a verification IR.} CHC-based encodings yield modular, solver-agnostic verification pipelines~\cite{bjorner2015horn, rybalchenko2011constraint}; Gurfinkel's notes~\cite{gurfinkel_chc_lecture} provide the linear CHC treatment directly applicable here.
    \item \textbf{Z3 Fixedpoint and SPACER.} The $\mu$Z engine~\cite{fixed-point-engine} provides CHC fixed-point computation within Z3; SPACER~\cite{z3guide_spacer} performs invariant synthesis and reachability queries, returning an inductive invariant or a counterexample.
    \item \textbf{Axiomatic semantics and SeaHorn.} Hoare-style semantics~\cite{utexas_axiomatic_semantics} underpin the state-transition encoding; SeaHorn~\cite{gurfinkel2015seahorn} demonstrates the same CHC architecture for C, which our pipeline adapts to Chiron's turtle-graphics semantics.
\end{itemize}

\subsection*{Implementation Details}

\subsubsection*{Extraction of symbolic program state}
\noindent\textit{\small File: \texttt{CHC\_Verification/variable\_name\_detection\_in\_IR.py}}

The verifier scans the linear Chiron IR to extract all user-declared variables and loop counters introduced by \texttt{repeat}, storing each with its corresponding Z3 \texttt{Real} variable. Internal repetition counters are kept separate from user variables to avoid name clashes; this determines the arity of the invariant predicate used throughout the pipeline.

\subsubsection*{Construction of the invariant relation}
\noindent\textit{\small File: \texttt{CHC\_Verification/z3\_fixed\_point.py}}

The mode-aware invariant predicate \texttt{Inv} captures the full program state: program counter, turtle position \texttt{(xcor, ycor)}, heading, pen state, and all user variables and loop counters. Four verification modes are supported:
\begin{enumerate}[leftmargin=1.5em]
    \item \textbf{default}: concrete start state, all values zero, \texttt{pendown = False};
    \item \textbf{universal}: turtle coordinates, heading, and user variables universally unconstrained;
    \item \textbf{specific}: user-supplied initial values, with unspecified variables defaulted to zero; and
    \item \textbf{safety-range}: ghost parameters recording initial user-variable values are appended to \texttt{Inv}, enabling the solver's invariant to be read as a precondition over starting values.
\end{enumerate}

\subsubsection*{Initialization of the CHC system}
\noindent\textit{\small File: \texttt{CHC\_Verification/init\_fixed\_point.py}}

The Z3 \texttt{Fixedpoint} object (engine: SPACER) is instantiated and seeded with a base-case constraint appropriate to the selected mode: a concrete \texttt{fp.fact} for \texttt{default}/\texttt{specific}, or a universally quantified \texttt{fp.rule} for \texttt{universal}/\texttt{safety-range}.

\subsubsection*{Translation from Chiron IR to CHC transition rules}
\noindent\textit{\small File: \texttt{CHC\_Verification/step\_rules.py}}

Each IR instruction is mapped to a universally quantified CHC implication over \texttt{Inv}:
\begin{itemize}[leftmargin=1.5em]
    \item \textbf{Instruction coverage:} assignments, conditional branching, assertions, \texttt{forward}/\texttt{backward} movement, \texttt{left}/\texttt{right} turns, \texttt{pendown}/\texttt{penup}, \texttt{goto}, \texttt{pause}, and \texttt{NOP}.
    \item \textbf{Expression translation:} arithmetic and boolean expressions are recursively translated to Z3, covering the four arithmetic operations, unary minus, all logical and comparison operators, boolean constants, and pen-status queries.
    \item \textbf{Turtle geometry:} \texttt{goto} yields exact coordinate updates; \texttt{forward}/\texttt{backward} use interval over-approximations of sine and cosine over discretized heading buckets; heading is normalized to $[0, 360)$ via a bounded sequence of conditional rewrites.
\end{itemize}

\subsubsection*{Property-checking interface}
\noindent\textit{\small File: \texttt{CHC\_Verification/safety\_properties.py}}

Given a user-provided boolean expression over the state variables, the verifier queries SPACER for a reachable violating state and reports \texttt{PASSED}, \texttt{FAILED}, or \texttt{UNKNOWN}. A counterexample is retained on failure; the synthesised invariant is stored on success. In \texttt{safety-range} mode, the invariant is additionally instantiated at the initial state to extract a symbolic safety precondition over starting values.

\subsection*{Usage}
The verifier is invoked from \texttt{Chiron-Framework/ChironCore/}. The four supported modes are:

\begin{verbatim}
uv run python CHC_Verification/safety_properties.py \
    <program.tl> <num_properties> default

uv run python CHC_Verification/safety_properties.py \
    <program.tl> <num_properties> universal

uv run python CHC_Verification/safety_properties.py \
    <program.tl> <num_properties> specific '{"x": 10, "y": 20}'

uv run python CHC_Verification/safety_properties.py \
    <program.tl> <num_properties> safety-range
\end{verbatim}

The tool then prompts interactively for property names and boolean expressions over \texttt{xcor}, \texttt{ycor}, \texttt{heading}, \texttt{pendown}, user variables, and loop counters; \texttt{And}, \texttt{Or}, and \texttt{Not} are also available.
